\section{Geometric Interpretation of the Integral}
\index{integral!geometric interpretation}%
\index{interpretation!geometric of the integral}%
%%%%

The following result establishes a correspondence between the Riemann integral
and the Peano-Jordan measure, providing a formal way to assign the Riemann
integral the meaning of the (signed) area of the subgraph of the integrand
function. In the following, we will never use this result, which serves only to
satisfy our intuition.

\begin{theorem}[geometric interpretation of the integral]
Let $f\colon[a,b]\to \RR$ be a bounded function, $a\leq b$. Define 
\begin{align*}
  E^+ &= \ENCLOSE{(x,y)\in [a,b]\times\RR\colon 0\le y \le f(x)}, \\
  E^- &= \ENCLOSE{(x,y)\in [a,b]\times\RR\colon f(x)\le y \le 0}   
\end{align*}
Then $f$ is Riemann-integrable on $[a,b]$ if and only if the sets $E^+$ and
$E^-$ are Peano-Jordan measurable. In such a case, it holds that:
\[
   \int_a^b f = m(E^+) - m(E^-)  
\]
where $m$ is the Peano-Jordan measure.
\end{theorem}
%
\begin{proof}
\emph{Step 1:} Suppose that $f(x)\ge 0$ for every $x\in [a,b]$. If $P=\ENCLOSE{x_0,x_1, \dots, x_N}$ is a partition of $[a,b]$, it is clear that, setting $I_k = [x_k,x_{k+1}]$, the union of the rectangles $R_k = I_k \times [0,\inf f(I_k)]$ gives us a polyrectangle contained in $E^+$. And $S_*(f,P)$ precisely represents the measure of such a polyrectangle. Thus, $m_*(E^+)\ge S_*(f,P)$ for every partition $P$, and therefore $m_*(E^+)\ge I_*(f,P)$. Reasoning analogously with the rectangles whose height is the $\sup$ of $f$, we observe that $S^*(f,P)$ represents the measure of a polyrectangle containing $E^+$, and thus $m^*(E^+)\le I^*(f,P)$.

On the other hand, given any polyrectangle contained in $E^+$, we can consider
the partition $P$ formed by $a$, $b$, and all the $x$ coordinates of the
vertical sides of the rectangles composing the polyrectangle. It is clear that
the polyrectangle will then be contained in the union of the rectangles
determined by $P$ with height the infimum of $f$. Thus, $m_*(E^+)\le I_*(f)$.
Similarly, if we take a polyrectangle that contains $E^+$, we can first cut all
the rectangles with the lines $x=a$ and $x=b$ and remove any rectangles that are
to the left of $x=a$ or to the right of $x=b$. This decreases the measure of the
polyrectangle and maintains the inclusion of $E^+$. At this point, consider the
partition $P$ obtained by taking the points $a$, $b$, and all the $x$
coordinates of the vertical sides of the rectangles that form the polyrectangle.
It is then clear that the measure of the original polyrectangle is greater than
or equal to $S^*(f,P)$, and thus $m^*(E^+)\ge I^*(f)$.

Therefore, since $m_*(E^+)=I_*(f)$ and $m^*(E^+)=I_*(f)$, it follows that if
$f\ge0$, then $f$ is Riemann-integrable if and only if $E^+$ is Peano-Jordan
measurable, and in such a case, we have $\int_a^b f = m(E^+)$.

\emph{Step 2:} Let $f\colon[a,b]\to \RR$ be arbitrary. Then we have $f = f^+ - f^-$. We can apply the previous step to $f^+$ and $f^-$. Observe that the set $E^+$ of $f^-$ is nothing but the reflection (with respect to the x-axis) of the set $E^-$ of $f$, and these sets have the same Peano-Jordan measure. Thus, if $f$ is Riemann-integrable, then $f^+$ and $f^-$ are also Riemann-integrable (Theorem~\ref{th:reticolo}), so $E^+$ and $E^-$ are Peano-Jordan measurable by the previous step, and we have 
\begin{equation}
  \label{47694}
\int_a^b f = \int_a^b f^+ - \int_a^b f^- = m(E^+) - m(E^-).  
\end{equation}
Conversely, if $E^+$ and $E^-$ are Peano-Jordan measurable, then $f^+$ and $f^-$
are Riemann-integrable (by the previous step), and also $f=f^+-f^-$ is
Riemann-integrable (Theorem~\ref{th:integrale_lineare}), and again \eqref{47694}
must hold.
\end{proof}

\begin{example}[area of the trapezoid]
Let $f(x) = mx + q$ be a linear function defined on an interval $[a,b]$, and
suppose that $f(a)\ge 0$ and $f(b)\ge 0$. The set $E^+=\ENCLOSE{(x,y)\colon x\in
[a,b], 0\le y \le f(x)}$ is thus a trapezoid with (vertical) bases $f(a)$ and
$f(b)$ and (horizontal) height $b-a$. Its area will thus be
\begin{align*}
  \int_a^b f 
  &= \int_a^b (mx+q)\, dx 
  = m\int_a^b x\, dx + \int_a^b q\, dx
  = \frac{m}{2}(b^2-a^2) + (b-a) q\\ 
  &= \frac 1 2 (b-a)(mb+ma+2q)
  = \frac 1 2 (b-a)(f(a) + f(b))
\end{align*}
which is indeed the elementary formula for calculating the area of a trapezoid.
\end{example}